% File: main.tex
\documentclass[letterpaper]{article}
\usepackage{aaai2026} % DO NOT CHANGE THIS
\usepackage{times}    % DO NOT CHANGE THIS
\usepackage{helvet}   % DO NOT CHANGE THIS
\usepackage{courier}  % DO NOT CHANGE THIS
\usepackage[hyphens]{url} % DO NOT CHANGE THIS
\usepackage{graphicx} % DO NOT CHANGE THIS
\urlstyle{rm}
\def\UrlFont{\rm}
\usepackage{natbib}   % DO NOT CHANGE THIS
\usepackage{caption}  % DO NOT CHANGE THIS
\frenchspacing        % DO NOT CHANGE THIS

\title{Security Threats and AI Defenses in 6G}

\author{
    Jodhveer Atwal\textsuperscript{\rm 1},
    Khushi Malik\textsuperscript{\rm 1},
    Mary Klassen\textsuperscript{\rm 1},
    Arda Yazici\textsuperscript{\rm 1},
    Victor Mai\textsuperscript{\rm 1}
}

\affiliations{
    \textsuperscript{\rm 1}Kwantlen Polytechnic University
}

\begin{document}
\maketitle

%%% -----------------------------------------------------
%%% SECTION: Introduction
%%% AUTHOR : Arda
%%% -----------------------------------------------------
\section{Introduction}
6G is expected to support highly dynamic, intelligent, and large-scale wireless systems that go beyond the capabilities of previous generations. Unlike earlier network generations, 6G is being designed with artificial intelligence integrated into core decision-making processes, including optimization, orchestration, and autonomous control. This shift creates opportunities for higher efficiency and adaptability, but it also introduces new security risks because attackers can target both network infrastructure and the learning systems that influence network behavior.

%%% -----------------------------------------------------
%%% SECTION: Background
%%% AUTHOR : Jodhveer
%%% -----------------------------------------------------
\section{Background}
6G networks are expected to combine ultra-low latency, distributed intelligence, massive device connectivity, and open architectures such as O-RAN. In this context, AI is used to support tasks such as spectrum management, anomaly detection, network slicing, and autonomous radio control.

%%% -----------------------------------------------------
%%% SECTION: Literature Review
%%% AUTHOR : Khushi, Mary, Victor
%%% -----------------------------------------------------
\section{Literature Review}

\subsection{Adversarial ML Threats and Defenses in AI-Enabled 6G}
% Khushi writes here.

\subsection{Broader 6G Adversarial Risk Surveys}
% Mary writes here.

\subsection{System-Level Threats, Multilayer Defense, and Secure Communication}
% Victor writes here.

%%% -----------------------------------------------------
%%% SECTION: Discussion
%%% AUTHOR : Victor
%%% -----------------------------------------------------
\section{Discussion}
The reviewed literature shows that AI improves the adaptability and automation of 6G systems, but also creates new security dependencies and new attack surfaces.

%%% -----------------------------------------------------
%%% SECTION: Proposed Idea
%%% AUTHOR : Jodhveer
%%% -----------------------------------------------------
\section{Proposed Idea}
A promising direction for future 6G security is a layered defense framework that combines trust-aware federated learning, adversarially robust model training, runtime anomaly monitoring, and explainable response policies.

%%% -----------------------------------------------------
%%% SECTION: Conclusion
%%% AUTHOR : Mary
%%% -----------------------------------------------------
\section{Conclusion}
The literature suggests that AI-driven 6G security remains an open problem and that future work must aim for more realistic and unified defense frameworks.

\bibliographystyle{aaai}
\bibliography{references}

\end{document}
